\documentclass[10pt,a4paper]{amsart}
\usepackage[margin=19mm]{geometry}
\usepackage{amsmath,amssymb,mathtools}
\usepackage[scr=rsfs,cal=euler]{mathalfa}
\usepackage{xfrac}
\usepackage[unicode,hidelinks,bookmarks=false]{hyperref}
\newcommand{\ie}{i.e.\ }
\newcommand{\cf}{cf.\ }
\newcommand{\resp}{respectively\ }
\newcommand{\Subsection}{\subsection}
\providecommand{\nobreakdash}{\nobreak\hbox{-}}
\setlength{\emergencystretch}{2em}
\multlinegap0pt
\usepackage{mathtools}
\usepackage[shortlabels]{enumitem}
\usepackage{amssymb}
\usepackage{xcolor}
\newtheorem{lemma}{Lemma}
\newcommand{\A}{\mathbb{A}} %Affine space
\newcommand{\C}{\mathcal{C}} %Cone
\newcommand{\DL}{\mathcal{D}} %Dirichlet--Lee domain
\newcommand{\Id}{\mathrm{Id}} %Identity map
\newcommand{\N}{\mathbb{N}} %Integers
\renewcommand{\P}{\mathbb{P}} %Projectivisation
\newcommand{\SA}{\mathrm{SA}} %Special Affine group
\newcommand{\st}{\; \middle\vert \;} % such that bar
\newcommand{\vp}[1]{{\left( {#1} \right)}} %between variable parentheses
\title{An Affine invariant Minkowski problem}
\author{Antoine Ablondi}
\date{Source: arXiv:2605.03791v1 (CC BY 4.0)}
\begin{document}
\maketitle

\noindent\emph{Source and licence.} An Affine invariant Minkowski problem. Version arXiv:2605.03791v1 (CC BY 4.0), distributed by arXiv under the Creative Commons Attribution 4.0 International licence, http://creativecommons.org/licenses/by/4.0/, as recorded in the arXiv licence metadata for this exact version and shown on the abstract page https://arxiv.org/abs/2605.03791v1. Full licence notice: \texttt{licenses/arxiv-2605.03791-CC-BY-4.0.txt}.\par\smallskip

\noindent\textit{Editorial scope.} Counted unit: the article's lemma domain (source0.tex lines 2278--2285). The article's complete original proof of it is delivered with it (source0.tex lines 2286--2318, one \texttt{proof} environment of 33 lines). No mathematics is written, repaired or re-derived: the statement and the proof are the article's own lines, in the article's own notation, and no retained line is substituted or re-flowed.  Retained uncounted as the source-local context the proof needs: lines 2230--2236.  Nothing was omitted from any retained span, so the package carries no omission record.  No retained line contains a citation, so the wrapper carries no bibliography and no bibliography entry.  No retained line contains a figure environment, an image-inclusion command or drawing code, so no image is delivered and the package declares no asset dependency.  Editorial formatting only, in addition: an AMS \texttt{amsart} preamble that restates the article's own declarations and macros used by the retained lines, namely the environments \texttt{lemma} on separate counters, together with the standard packages those lines need.  The retained environments print in this document as Lemma 1 in order of appearance, whereas the article prints the same items with the numbering of its own full document.  The complete native source of the article as distributed in the arXiv e-print for this exact version is bundled unchanged as \texttt{sources/199815/source0.tex}.
\begin{lemma}
\label{lem scalar product to - infty}
Let $[\varphi] \in \P(\C^*)$ and $A$ be a compact subset of $D_\tau$. If $(\gamma^n_\tau)_{n\in \N}$ is a diverging sequence of elements of $\Gamma_\tau$, then for all $P \in \A^{d+1}$,
\begin{equation*}
 \sup_{Q \in A} \varphi\vp{ \overrightarrow{P(\gamma^n_\tau Q)}} \xrightarrow[n \to +\infty]{} -\infty \, . 
\end{equation*}
\end{lemma}

\begin{lemma}
\label{lem int of Dirichlet--Lee domain}
Let $[\varphi]\in \P(\C^*)$. The Dirichlet--Lee domain $\DL_{[\varphi]}$ of $D_\tau$ has interior
\begin{equation*}
 \mathrm{int}(\DL_{[\varphi]}) = \left\lbrace P \in D_\tau \st \forall \gamma_\tau \in \Gamma_\tau\setminus \{\Id\}, \ \varphi\vp{\overrightarrow{P(\gamma_\tau P)}} < 0\right\rbrace ,
\end{equation*}
where $\Id$ is the identity in $\SA (\A^{d+1})$. 
\end{lemma}

\begin{proof}
The direct inclusion $\subseteq$ is trivial. Let us then prove the reverse inclusion.

Let us work in the parametrisation $(\mathbf{P})$. We now let $Y=(y,-1) \in \C^*$ and we want to show that 
\begin{equation*}
\left\lbrace X \in D_\tau \st \forall (\gamma,\tau) \in \Gamma_\tau\setminus \{(I,0)\}, \ X\cdot Y > (\gamma X + \tau ) \cdot Y\right\rbrace \subseteq \, \mathrm{int}(\DL_{y}) \,.
\end{equation*}
Let $X \in D_\tau$ such that for all $(\gamma,\tau) \in \Gamma_\tau\setminus \{(I,0)\}$, we have $X\cdot Y > (\gamma X + \tau ) \cdot Y$. Let $A$ be a compact neighbourhood of $X$ in $D_\tau$, such that 
\begin{equation}
\label{eq1 lem int}
\inf_{X' \in A} X' \cdot Y \geq X \cdot Y -1 \, . 
\end{equation}
Lemma~\ref{lem scalar product to - infty} and the proper discontinuity of the action of $\Gamma_\tau$ on $D_\tau$ imply that there is only a finite family of elements $(\gamma_i,\tau_i)_{1 \leq i \leq n}$ of $\Gamma_\tau \setminus\{(I,0)\}$ such that for all $(\gamma,\tau) \in \Gamma_\tau \setminus\{(I,0),(\gamma_1,\tau_1), \dots, (\gamma_n,\tau_n)\}$,
\begin{equation}
\label{eq2 lem int}
X \cdot Y - 1 > \sup_{X' \in A} (\gamma X' + \tau) \cdot Y \, . 
\end{equation}
Thus, using equations~\eqref{eq1 lem int} and \eqref{eq2 lem int}, for all $X' \in A$ and $(\gamma,\tau) \in \Gamma_\tau \setminus\{(I,0),(\gamma_1,\tau_1), \dots, (\gamma_n,\tau_n)\}$,
\begin{equation}
\label{eq3 lem int}
 X' \cdot Y > (\gamma X' + \tau ) \cdot Y \, . 
\end{equation}
Now, for all $1 \leq i \leq n$, we have, by assumption, that $X\cdot Y > (\gamma_i X + \tau_i ) \cdot Y$, so there is a neighbourhood $U_i \subseteq A$ of $X$ such that for all $X' \in U_i$,
\begin{equation}
\label{eq4 lem int}
X' \cdot Y >(\gamma_i X' + \tau_\gamma) \cdot Y \, . 
\end{equation}
Equations~\eqref{eq3 lem int} and \eqref{eq4 lem int} imply the neighbourhood $U \coloneqq \bigcap_{1 \leq i \leq n} U_i \subseteq A$ of $X$ in $D_\tau$ satisfies that for all $X' \in U$ and $(\gamma,\tau) \in \Gamma_\tau$, 
\begin{equation*}
 X' \cdot Y \geq (\gamma X' + \tau ) \cdot Y \, . 
\end{equation*}
That means that $X$ admits a neighbourhood included in $\DL_{y}$, i.e. $X \in \mathrm{int}(\DL_{y})$. 
\end{proof}

\end{document}
